% Position markers for the items of [M] (yang-mills-informalization, manuscript/main.tex), \input
% before \begin{document}.  Adapted from the ternary Goldbach companion's regionmarks.tex; the
% marker mechanism (zref-savepos labels in a list of their own, written at shipout) is
% formal-companion's informalize fcmarkers.sty \fc@savepos.  Every marker is a zero-size whatsit
% placed where it adds no break point and changes no \lastskip; layoutcheck.py compares the
% rebuilt PDF with the committed one word by word.
%
%   thm:<n>:head   theorem-like head, inside the label box: the baseline of the first line
%   thm:<n>:end    the end of the last line (horizontal mode only), and
%   thm:<n>:below  after that line on the vertical list (also: after a closing list or paragraph)
%   \rmenddisp{n}{j} in the .aux: theorem n ends with display j (no node is added)
%   list:<l>:below after the last line of the l-th list in a theorem-like item, and
%   \rmendlist{n}{l} in the .aux: theorem n ends with list l
%   dend:<j>       after the box of the j-th display (equation, equation*, \[ \], and the
%                  amsmath alignments, whose last row is followed by a \noalign)
%   eq:<k>:pre     after the line before the k-th numbered display environment, whose last space
%                  it removes as the display would (horizontal mode,
%                  and not on the empty list after a display, where a \vadjust would add a line)
%   \rmafterdisp{k} in the .aux: the k-th numbered environment opens such an empty list
%   eq:<k>:tag:<r> the r-th number of the k-th numbered environment, left of "(" on its baseline
%   eq:<k>:post    after the box of the k-th numbered environment (not for multline)
\RequirePackage{zref-savepos}
\RequirePackage{zref-abspage}
\makeatletter
\zref@newlist{rmpos}
\zref@addprops{rmpos}{posx,posy,abspage}
\newcommand*\rm@savepos[1]{%
  \begingroup
    \edef\rm@tmp{\endgroup
      \noexpand\ZREF@zsavepos\noexpand\zref@labelbylist{rmpos}{#1}}%
  \rm@tmp}
\AddToHook{begindocument/end}{%
  \immediate\write\@auxout{\string\rmgeometry{pdfpagewidth=\number\pdfpagewidth,%
    pdfpageheight=\number\pdfpageheight,%
    textleft=\number\dimexpr1in+\hoffset+\oddsidemargin\relax,%
    evenleft=\number\dimexpr1in+\hoffset+\evensidemargin\relax,%
    texttop=\number\dimexpr1in+\voffset+\topmargin+\headheight+\headsep\relax,%
    textwidth=\number\textwidth,textheight=\number\textheight,%
    baselineskip=\number\baselineskip}}}
\providecommand\rmgeometry[1]{}
% display ends
\newcount\rm@dk
\def\rm@dend{\global\advance\rm@dk\@ne\rm@savepos{dend:\the\rm@dk}}
\newif\ifrm@num
\newcount\rm@eqk
\AddToHook{env/equation/end}{\vadjust{\rm@dend\rm@savepos{eq:\the\rm@eqk:post}}}
\AddToHook{env/equation*/end}{\vadjust{\rm@dend}}
\let\rm@orig@black\black@
% (split, inside an equation, also ends with \black@: no marker there)
\newif\ifrm@split
\AddToHook{env/split/begin}{\global\rm@splittrue}
\AddToHook{env/split/after}{\global\rm@splitfalse}
\def\black@#1{\rm@orig@black{#1}\ifmeasuring@\else\ifrm@split\else
  \noalign{\rm@dend\ifrm@num\rm@savepos{eq:\the\rm@eqk:post}\fi}\fi\fi}
% theorem-like items
\gdef\rm@thm{}
\let\rm@orig@thmhead\thmhead
\def\thmhead#1#2#3{\xdef\rm@thm{#2}\global\rm@inthmtrue\global\rm@listlastfalse\rm@savepos{thm:#2:head}\rm@orig@thmhead{#1}{#2}{#3}}
\def\rm@thmend{%
  \ifhmode
    \ifnum\lastnodetype=\m@ne
      \rm@enddisp
    \else
      \unskip\rm@savepos{thm:\rm@thm:end}\vadjust{\rm@savepos{thm:\rm@thm:below}}%
    \fi
  \else
    \ifnum\lastnodetype=\@ne \rm@savepos{thm:\rm@thm:below}\else
    \ifnum\lastnodetype=\tw@ \rm@savepos{thm:\rm@thm:below}\else
    \ifrm@listlast \rm@endlist\else
      \rm@enddisp
    \fi\fi\fi
  \fi\global\rm@inthmfalse}
% a list inside a theorem-like item: the end of its last item's last line (lists end in vertical
% mode, on glue, where a node would change \lastskip)
\newif\ifrm@inthm
\newif\ifrm@listlast
\newcount\rm@lk
\def\rm@listend{\ifrm@inthm\ifhmode\global\advance\rm@lk\@ne
  \unskip\vadjust{\rm@savepos{list:\the\rm@lk:below}}\global\rm@listlasttrue\fi\fi}
\def\rm@endlist{\if@filesw\immediate\write\@auxout{\string\rmendlist{\rm@thm}{\the\rm@lk}}\fi}
\providecommand\rmendlist[2]{}
\AddToHook{env/itemize/end}{\rm@listend}
\AddToHook{env/enumerate/end}{\rm@listend}
\AddToHook{env/itemize/after}{\ifrm@inthm\global\rm@listlasttrue\fi}
\AddToHook{env/enumerate/after}{\ifrm@inthm\global\rm@listlasttrue\fi}
\AddToHook{para/begin}{\global\rm@listlastfalse}
% a theorem that ends with a display: name that display's end marker, without a node
\def\rm@enddisp{\if@filesw\immediate\write\@auxout{\string\rmenddisp{\rm@thm}{\the\rm@dk}}\fi}
\providecommand\rmenddisp[2]{}
\AddToHook{env/maintheorem/end}{\rm@thmend}
\AddToHook{env/theorem/end}{\rm@thmend}
\AddToHook{env/proposition/end}{\rm@thmend}
\AddToHook{env/lemma/end}{\rm@thmend}
\AddToHook{env/corollary/end}{\rm@thmend}
\AddToHook{env/definition/end}{\rm@thmend}
\AddToHook{env/remark/end}{\rm@thmend}
\AddToHook{env/principal/end}{\rm@thmend}
% numbered display environments
\def\rm@numbegin{\global\advance\rm@eqk\@ne\global\rm@tagr\z@\global\rm@indisptrue
  \ifhmode\ifnum\lastnodetype=\m@ne\rm@afterdisp\else\unskip\vadjust{\rm@savepos{eq:\the\rm@eqk:pre}}\fi\fi}
% a display that opens an empty list (it follows a display directly): say so, without a node
\def\rm@afterdisp{\if@filesw\immediate\write\@auxout{\string\rmafterdisp{\the\rm@eqk}}\fi}
\providecommand\rmafterdisp[1]{}
\AddToHook{env/equation/begin}{\rm@numbegin}
\AddToHook{env/align/begin}{\rm@numbegin\global\rm@numtrue}
\AddToHook{env/align/after}{\global\rm@numfalse\global\rm@indispfalse}
\AddToHook{env/multline/begin}{\rm@numbegin}
\newcount\rm@tagr
% numbers outside a numbered environment (\eqref) get no marker
\newif\ifrm@indisp
\let\rm@orig@tagform\tagform@
\def\tagform@#1{\ifrm@indisp\ifmeasuring@\else\global\advance\rm@tagr\@ne\fi
  \maketag@@@{\rm@savepos{eq:\the\rm@eqk:tag:\the\rm@tagr}(\ignorespaces#1\unskip\@@italiccorr)}%
  \else\rm@orig@tagform{#1}\fi}
\AddToHook{env/equation/after}{\global\rm@indispfalse}
\AddToHook{env/multline/after}{\global\rm@indispfalse}
\makeatother
